% Generated draft: requires mathematical and source review.
\documentclass[11pt]{article}
\usepackage[margin=25mm]{geometry}
\usepackage{fontspec}
\setmainfont{Latin Modern Roman}
\usepackage{amsmath,amssymb,mathtools,unicode-math}
\setmathfont{Latin Modern Math}
\usepackage{enumitem,needspace,xurl,hyperref}
\setlength{\parindent}{0pt}
\setlength{\parskip}{5pt}
\setlength{\emergencystretch}{3em}
\newcommand{\N}{\mathbb{N}}
\newcommand{\Z}{\mathbb{Z}}
\newcommand{\Q}{\mathbb{Q}}
\newcommand{\R}{\mathbb{R}}
\newcommand{\C}{\mathbb{C}}
\newcommand{\F}{\mathbb{F}}
\newcommand{\A}{\mathbb{A}}
\newcommand{\PP}{\mathbb{P}}
\newcommand{\E}{\mathbb{E}}
\newcommand{\eps}{\varepsilon}
\newcommand{\dd}{\,\mathrm{d}}
\begin{document}
\section*{UnsolvedMath ID: 30000129}
\textbf{Source problem code:} catalogue entry
\subsection{Definitions and mathematical statement}
The source asks whether, in the intended hierarchical mean-field model, the continuum limit and the limit \(N\to\infty\) may be interchanged. It does not specify the stochastic model, what \(N\) counts, the continuum scaling, initial laws, time scales, or the mode of convergence. Consequently the following is a notation scheme for the question, not a theorem for arbitrary hierarchical systems.

Write \(\varepsilon\downarrow0\) for the unspecified continuum approximation and \(X^{N,\varepsilon}\) for the corresponding stochastic objects, with \(N\in\Z_{>0}\). If the intended convergence is weak convergence of laws, one must first identify these objects in a common state or path space \(E\), using the appropriate rescalings and embeddings, and set \(\mu_{N,\varepsilon}=\mathcal L(X^{N,\varepsilon})\). The interchange question would then ask for existence and equality, in the intended topology on laws on \(E\), of
\[
 \lim_{\varepsilon\downarrow0}\left(\lim_{N\to\infty}\mu_{N,\varepsilon}\right)
 \quad\text{and}\quad
 \lim_{N\to\infty}\left(\lim_{\varepsilon\downarrow0}\mu_{N,\varepsilon}\right).
\]
Both inner limits and both outer limits must be defined in compatible spaces. If the original objects are generators, semigroups, correlation functions or renormalization maps instead of process laws, this display must be replaced by the corresponding model-specific formulation; the record does not select one. Nor does it state whether interchange is sought for every admissible initial state, uniformly in time, or only for particular observables. Equality of two iterated limits is a different target from a joint limit along every sequence \((N,\varepsilon)\to(\infty,0)\); no joint-limit assertion is added.

The literature triage checked on 2026-08-21 labels the problem partially solved. It reports implementations of the interchange and renormalization program for important hierarchical branching and diffusion models, including two-type branching models. It expressly qualifies the conclusions as model-dependent and states that the central two-type paper remains conditional on a boundary-behavior conjecture. Neither that conjecture nor the applicable hypotheses are supplied. These reports do not establish the interchange for an unspecified general class of models, and their current status has not been independently checked here.
\subsection{Short English statement}
The question asks whether taking a continuum approximation and taking the hierarchical mean-field limit in either order gives the same limiting object. A precise assertion needs the underlying model, scalings, initial data and topology. The supplied partial progress is model-dependent and retains an explicit boundary-behavior conjecture in the central two-type case.
\subsection{Sources}
\begin{itemize}
\item Local UnsolvedMath registry, permanent ID 30000129, source code catalogue entry, ``Interchanging Hierarchical Mean-Field and Continuum Limits,'' including the original statement, cleaned statement and supplied verification and literature fields.
\item \emph{Large Scale Stochastic Dynamics (2005)}, Oberwolfach Reports, pp.\ 2229--2310, DOI 10.4171/owr/2004/43: \url{https://doi.org/10.4171/owr/2004/43}. The supplied report title uses 2005, whereas the DOI and proposed-year metadata use 2004.
\item Supplied reference ``Dawson--Greven--den Hollander--Sun--Swart, two-type branching'': \url{https://arxiv.org/abs/math/0610645}.
\item Supplied reference ``Published renormalization theorem'': \url{https://doi.org/10.1214/07-AIHP143}.
\item Literature triage checked 2026-08-21; registry record updated 2026-09-16. The record reports official-source verification and normalized wording; this draft has not visited the cited sources or independently verified those claims.
\end{itemize}
\end{document}
